\documentclass[compact]{ne630handout}
\HandoutTypography{9}{10.7}{5pt}

\handoutsolutionsfalse
% Uncomment to reveal numerical answers and instructor cues.
% \handoutsolutionstrue
\ifdefined\instructorsolutions\handoutsolutionstrue\fi

\hypersetup{pdftitle={NE 630 -- Lesson 28: Precursor Groups and the Inhour Equation},
  pdfauthor={Jeremy Roberts}}

\begin{document}
\Lesson{28}{Precursor Groups and the Inhour Equation}{FNRP 5.4--5.5, especially Table 5.1.}
{Construct a grouped-precursor model, derive the inhour equation,
and compare reactor periods from one- and six-group models.}

\begin{handoutflow}
\Sec{1}{From one nuclide to precursor groups}

Many nuclides produce delayed neutrons. For thermal fission of U-235,
the ten examples below account for about 72\% of delayed neutrons
in the cited evaluation.
The last column gives each nuclide's percentage of the total delayed-neutron
yield, not a percentage of all fission neutrons.

% Shares: Dimitriou et al. (2021), Table VII, U-235 thermal, CRP column.
% Half-lives: Tables I--II; As-85 from the linked ENSDF evaluation.
\begin{center}
{\fontsize{8}{9.5}\selectfont
\begin{tabular}{lrr}
\toprule
\textbf{Precursor} & $t_{1/2}$ (s) & \textbf{Share of delayed yield (\%)}\\
\midrule
$^{137}$I  & 24.6  & 16.3\\
$^{89}$Br  & 4.34  & 11.1\\
$^{94}$Rb  & 2.70  & 9.3\\
$^{90}$Br  & 1.91  & 7.4\\
$^{88}$Br  & 16.3  & 7.3\\
$^{85}$As  & 2.02  & 5.3\\
$^{138}$I  & 6.25  & 4.7\\
$^{139}$I  & 2.28  & 3.5\\
$^{95}$Rb  & 0.378 & 3.5\\
$^{87}$Br  & 55.6  & 3.2\\
\midrule
\textbf{Subtotal} & & \textbf{71.6}\\
\bottomrule
\end{tabular}}
\end{center}
\SourceLine{Yield shares: P. Dimitriou et al.,
\href{https://doi.org/10.1016/j.nds.2021.04.006}
{\emph{Nuclear Data Sheets} 173, 144--238 (2021)}, Table VII,
U-235 thermal, CRP column (CRP+ENDF/B-VIII.0 decay data;
JEFF-3.1.1 fission yields). Half-lives rounded from Tables I--II;
As-85 from \href{https://www.nndc.bnl.gov/ensnds/85/Se/beta_decay_2.021_s.pdf}
{B. Singh and J. Chen, ENSDF (2014)}.}

\textbf{Average by delayed-neutron yield.} Let $s_j$ be the delayed-yield
share of nuclide $j$. Normalize the shares over the precursors being averaged:
\[
 \boxed{\bar{t}_{1/2}=\frac{\sum_j s_j t_{1/2,j}}{\sum_j s_j},\qquad
 \bar{\lambda}=\frac{\ln2}{\bar{t}_{1/2}}.}
\]
This preserves the mean delay $\bar{t}_{1/2}/\ln2$.
For the ten listed nuclides, divide by $71.6$: $\bar{t}_{1/2}=11.66$ s
and $\bar{\lambda}=0.05946$ s$^{-1}$ \emph{for this subset}.
The remaining $28.4\%$ is needed for a full U-235 average.
\SourceLine{Averaging procedure: FNRP Eqs. (5.32)--(5.34).}

\columnbreak
Represent the combined release of all precursors
with $G$ effective groups, each with delayed fraction $\beta_i$ and
decay constant $\lambda_i=\ln2/t_{1/2,i}$.
The values below are from FNRP~Table~5.1.

\begin{center}
{\fontsize{8}{9.5}\selectfont
\setlength{\tabcolsep}{4pt}
\begin{tabular}{crrrr}
\toprule
 & & \multicolumn{3}{c}{\textbf{Delayed fraction $\beta_i$}}\\
\cmidrule(lr){3-5}
$i$ & $t_{1/2,i}$ (s) & U-233 & U-235 & Pu-239\\
\midrule
1 & 56   & 0.00023 & 0.00021 & 0.00007\\
2 & 23   & 0.00078 & 0.00142 & 0.00063\\
3 & 6.2  & 0.00064 & 0.00128 & 0.00044\\
4 & 2.3  & 0.00074 & 0.00257 & 0.00069\\
5 & 0.61 & 0.00014 & 0.00075 & 0.00018\\
6 & 0.23 & 0.00008 & 0.00027 & 0.00009\\
\midrule
\multicolumn{2}{l}{$\beta=\sum_i\beta_i$}
  & 0.00261 & 0.00650 & 0.00210\\
\multicolumn{2}{l}{$\bar\nu$ (neutrons/fission)}
  & 2.50 & 2.43 & 2.90\\
\bottomrule
\end{tabular}}
\end{center}
Each $C_i$ is an \textbf{effective delayed-neutron reservoir}.
Unlike Lesson 27's $C=f_nN_{\rm Br}$, a grouped $C_i$ alone does not
determine an individual nuclide's concentration.

\Sec{2}{Equations and initial conditions}

With $\beta=\sum_{i=1}^{G}\beta_i$,
\[
 \boxed{\begin{aligned}
 \frac{dn}{dt}&=\frac{\rho(t)-\beta}{\Lambda}n+
                 \sum_{i=1}^{G}\lambda_i C_i+S(t),\\
 \frac{dC_i}{dt}&=\frac{\beta_i}{\Lambda}n-\lambda_i C_i,
                       \qquad i=1,\ldots,G .
 \end{aligned}}
\]
For an initially critical, source-free equilibrium,
\[
 \rho(0^-)=0,\qquad n(0)=n_0,\qquad
 \boxed{C_i(0)=\frac{\beta_i n_0}{\lambda_i\Lambda}.}
\]
These equations have the matrix form
$\dot{\mathbf y}=\mathbf A(\rho)\mathbf y+\mathbf b$,
where $\mathbf y=[n,C_1,\ldots,C_G]^T$ and
$\mathbf b=[S,0,\ldots,0]^T$.
The $7\times7$ matrix for six groups has
\[
 \begin{aligned}
 A_{00}&=(\rho-\beta)/\Lambda,& A_{0i}&=\lambda_i,\\
 A_{i0}&=\beta_i/\Lambda,& A_{ii}&=-\lambda_i
 \end{aligned}
\]
and all other entries zero.

\textbf{Check.} With U-235, $n_0=10^8$ cm$^{-3}$, and
$\Lambda=5\times10^{-5}$ s, calculate $C_1(0)$.
Why is setting all $C_i(0)=0$ a different physical experiment?
\[
 C_1(0)=\RevealMath[1.4in]{3.393\times10^{10}}\ {\rm cm}^{-3}.
\]

\end{handoutflow}

\clearpage
\HandoutContinue{Reactor periods and a numerical comparison}
\begin{handoutflow}
\Sec{3}{The pathway to the inhour equation}

After a step, take $\rho$ constant and $S=0$.
\textbf{Try one exponential mode:}
\[
 n(t)=a e^{\omega t},\qquad C_i(t)=b_i e^{\omega t}.
\]
\textbf{Eliminate the precursor amplitudes:}
\[
 (\omega+\lambda_i)b_i=\frac{\beta_i}{\Lambda}a
 \quad\Longrightarrow\quad
 b_i=\frac{\beta_i a}{\Lambda(\omega+\lambda_i)}.
\]
\textbf{Substitute into the neutron equation:}
\[
 \Lambda\omega=\rho-\beta+
                  \sum_{i=1}^{G}\frac{\beta_i\lambda_i}{\omega+\lambda_i}.
\]
\textbf{Use $\beta=\sum_i\beta_i$ and rearrange:}
\[
 \boxed{\rho=\Lambda\omega+
          \sum_{i=1}^{G}\frac{\beta_i\omega}{\omega+\lambda_i}.}
\]
This is the \textbf{inhour equation}. It is rational; clearing the
denominators gives the degree-$(G+1)$ characteristic polynomial.
The poles $-\lambda_i$ are not roots.

The full solution is a sum of $G+1$ modes, with amplitudes fixed by the
initial state. The largest eigenvalue $\omega_\star$ determines the
late-time signed period $T=1/\omega_\star$:
\[
 \boxed{\rho=\frac{\Lambda}{T}
            +\sum_{i=1}^{G}\frac{\beta_i}{1+\lambda_iT}.}
\]
At criticality, $\omega_\star=0$ and the equilibrium period is infinite.
For a negative insertion, $T<0$.

\Sec{4}{One effective group and useful limits}

\textbf{Near criticality}, if $|\omega_\star|\ll\lambda_i$ for every group,
\[
 T\simeq\frac{\Lambda+\sum_i\beta_i/\lambda_i}{\rho}.
\]
A one-group model using Section 1's average preserves the equilibrium
total precursor inventory:
\[
 \beta_{\rm 1g}=\beta,\qquad
 \lambda_{\rm 1g}=\bar\lambda=\frac{\beta}{\sum_i\beta_i/\lambda_i}.
\]
For the U-235 groups on page 1, $\bar t_{1/2}=9.044$ s and
$\bar\lambda=0.076641$ s$^{-1}$. This reproduces the small-reactivity
limit. A mixture still has several modes; one effective group cannot
reproduce every transient. This is distinct from L27's Br-87-only model.

\textbf{Well above prompt criticality},
\[
 T\simeq\frac{\Lambda}{\rho-\beta}.
\]
This estimate requires the delayed term to be small relative to
$\rho-\beta$; the full period is finite at $\rho=\beta$.

\textbf{For a sufficiently large negative insertion},
\[
 \omega_\star\longrightarrow-\lambda_1,\qquad
 T\longrightarrow-\frac{1}{\lambda_1}=-80.79\ {\rm s}.
\]
Earlier decay can be much faster than this late-time limit.

\columnbreak
\Sec{5}{Work with \texttt{ne630\_kinetics}}

Save \texttt{ne630\_kinetics.py} beside your notebook or script.
Use Python with NumPy and SciPy; add Matplotlib for plots.
Import it as \texttt{import ne630\_kinetics as kin}.
The defaults are the six U-235 groups, $\Lambda=5\times10^{-5}$ s,
critical initial equilibrium, $S=0$, and no feedback.
Reactivity arguments are dimensionless: $\rho=d\,\texttt{kin.BETA}$
for an insertion of $d$ dollars. Times are in seconds.

\textbf{From a terminal}, run
\mbox{\texttt{python ne630\_kinetics.py}} to print period tables and
checks; it also writes \texttt{numerical\_results.json} beside the module.
Import the module in your own notebook, script, or interactive Python
session to vary the cases and make plots.

\textbf{For a step}, use \texttt{kin.sample\_times(tmax)} and
\texttt{kin.solve\_step(rho, times)}. The result \texttt{sol} contains
times in \texttt{sol.t}, $n/n_0$ in \texttt{sol.y[0]}, and $C_i/n_0$ in
\texttt{sol.y[i]} ($i=1,\ldots,6$).
\texttt{kin.dominant\_period(rho)} returns the signed period.
For one group, pass \texttt{beta=[kin.BETA]} and
\texttt{lam=[lambda\_eff]} to both functions, with
\texttt{lambda\_eff} set to the full six-group $\bar\lambda$ from Section 4.
Each model initializes its own critical precursor inventory.

\Sec{6}{Compare one group with six}

\textbf{Verify equilibrium.} Set $\rho=0$: all densities should remain
constant. Inspect \texttt{kin.matrix} to connect the implementation to
Section 2.

\textbf{Compute periods.} Run $\rho/\beta=0.02$ and $0.2$.
Compare the six-group period, the effective one-group period, and the
small-reactivity estimate in Section 4.
\begin{center}
\renewcommand{\arraystretch}{1.45}
\begin{tabular}{@{}lcc@{}}
\toprule
$\rho/\beta$ & $T_{\rm 6g}$ (s) & $T_{\rm 1g}$ (s)\\
\midrule
0.02 & \RevealMath[.82in]{617.6} & \RevealMath[.82in]{639.7}\\
0.2  & \RevealMath[.82in]{36.88} & \RevealMath[.82in]{52.24}\\
\bottomrule
\end{tabular}
\end{center}
Use \texttt{kin.inhour\_residual} to check that each six-group
frequency $1/T$ satisfies the inhour equation at the imposed $\rho$.

\textbf{Plot and explain.} Plot both $0.2\beta$ traces over
$0$--$200$ s on semilog axes. Compare the late slope with $1/T$.
Why does preserving the mean delay and equilibrium inventory fail to
preserve the whole transient? Which insertion is better represented
by one effective group?

\textbf{Notebook enrichment.} Continue in
\nolinkurl{lesson_28_exploration.ipynb} for the closer-to-critical
$0.002\beta$ case, prompt-critical crossings, negative-reactivity tails,
and finite pulses that demonstrate precursor memory.
The notebook uses the same module and six-group data.
\end{handoutflow}
\end{document}
